\documentclass{article} % For LaTeX2e
\usepackage{iclr2027_conference,times}

% Optional math commands from https://github.com/goodfeli/dlbook_notation.
\input{math_commands.tex}

\usepackage{hyperref}
\usepackage{url}
\usepackage{graphicx}


\title{Dense Semantic Priors for Diffusion-Based BEV Representation Learning}

% Authors must not appear in the submitted version. They should be hidden
% as long as the \iclrfinalcopy macro remains commented out below.
% Non-anonymous submissions will be rejected without review.

\author{Hyelin Jo \thanks{ Use footnote for providing further information
about author (webpage, alternative address)---\emph{not} for acknowledging
funding agencies.  Funding acknowledgements go at the end of the paper.} \\
Automotive Intelligence X Electrification Lab (AXELab) \\
Cho Chun Shik Graduate School of Mobility, KAIST \\
Daejeon, South Korea \\
\texttt{hyelin.jo@kaist.ac.kr} \\
\And
Dongseok Kum \\
Automotive Intelligence X Electrification Lab (AXELab) \\
Cho Chun Shik Graduate School of Mobility, KAIST \\
Daejeon, South Korea \\
\texttt{professor_email@kaist.ac.kr}
}


% The \author macro works with any number of authors. There are two commands
% used to separate the names and addresses of multiple authors: \And and \AND.
%
% Using \And between authors leaves it to \LaTeX{} to determine where to break
% the lines. Using \AND forces a linebreak at that point. So, if \LaTeX{}
% puts 3 of 4 authors names on the first line, and the last on the second
% line, try using \AND instead of \And before the third author name.

\newcommand{\fix}{\marginpar{FIX}}
\newcommand{\new}{\marginpar{NEW}}

%\iclrfinalcopy % Uncomment for camera-ready version, but NOT for submission.
\begin{document}


\maketitle

\begin{abstract}
Bird's-eye-view (BEV) representations are central to camera-based autonomous
driving perception, yet they are typically learned only through downstream task
supervision. Recent BEV diffusion methods provide feature-level denoising
objectives, but often rely on task-specific annotations or conditioning signals,
which limits their applicability across perception tasks. We propose a
semantic-guided diffusion framework for BEV representation learning that uses
dense semantic priors instead of task-specific conditioning. Multi-view semantic
masks extracted from a frozen vision foundation model are projected into the BEV
space to construct a spatially aligned semantic BEV prior. This prior guides BEV
feature denoising through Semantic BEV Adaptive Modulation (SBAM), enabling
generative diffusion pre-training without downstream task supervision. The
pre-trained diffusion model is then used to provide feature-level supervision
for BEV perception models through masked generative distillation. At inference
time, both the diffusion model and semantic branch are removed. The proposed
framework achieves competitive performance against task-specific diffusion
approaches without additional inference cost.
\end{abstract}

%% Please note that we have introduced automatic line number generation
%% into the style file for \LaTeXe. This is to help reviewers
%% refer to specific lines of the paper when they make their comments. Please do
%% NOT refer to these line numbers in your paper as they will be removed from the
%% style file for the final version of accepted papers.

\section{Introduction}
Bird's-eye-view (BEV) representations have become fundamental to autonomous
driving perception by integrating multi-view images into a unified top-down
space~\citep{SUR1,SUR2,BVF}. They are widely used for diverse downstream tasks,
including 3D object detection, BEV segmentation, motion prediction, and
end-to-end planning~\citep{SUR2,BVF,LSS,UniAD,VAD}. However, camera-based BEV
features are typically learned indirectly through downstream task losses without
explicit supervision in the latent feature space. Consequently, errors and
uncertainties introduced during image-to-BEV projection may propagate into the
resulting BEV representations~\citep{BEVDiff,DiffBEV}.

Recent studies have introduced diffusion models into BEV perception to improve
BEV representation quality~\citep{DiffBEV,DifFUSER,BEVDiff}. Diffusion models
learn the underlying data distribution by progressively adding noise to samples
and training a network to reverse the corruption process~\citep{DDPM,DDIM}.
When applied to BEV features, the resulting denoising objective provides an
explicit feature-level training signal for recovering spatially structured
representations.

Despite their potential, existing BEV diffusion methods often rely on
task-specific priors or ground-truth information for
conditioning~\citep{BEVDiff,LaneDiff}. For example, BEVDiffuser constructs a
ground-truth object layout from object categories and 3D bounding
boxes~\citep{BEVDiff}, whereas LaneDiffusion uses lane centerline
priors~\citep{LaneDiff}; both employ conditioning mechanisms tailored to their
respective representations. Although such conditioning provides effective
guidance, it requires task-specific annotations and ties the diffusion model to
particular perception tasks. As a result, adapting these models to other tasks
may involve redesigning both the conditioning signal and its injection
mechanism.

Task dependency also appears in the training objective. Existing BEV diffusion
methods commonly optimize the diffusion model together with downstream task
losses~\citep{BEVDiff,DiffBEV,DifFUSER,DiffVecMap}, directly coupling diffusion
training to specialized prediction heads and supervision. Although
BEVDiffuser~\citep{BEVDiff} removes its diffusion module at inference time, its
diffusion training still relies on a 3D detection objective. Extending these
approaches to another perception task may therefore require modifying the
prediction head and training objectives, limiting the applicability of
diffusion-based BEV representation learning.

In this work, we propose a task-agnostic semantic-guided diffusion framework for
BEV representation learning, as illustrated in Figure~\ref{fig:overview}. The
framework is task-agnostic in both its conditioning design and training
objective, avoiding task-specific annotations and downstream task supervision
during diffusion pre-training. A frozen vision foundation model~\citep{SAM3}
extracts multi-view semantic masks, which are projected into the BEV space to
construct a spatially aligned dense semantic BEV prior. This prior captures both
dynamic agents and static scene elements and guides denoising through Semantic
BEV Adaptive Modulation (SBAM). After pre-training, the diffusion model is
frozen, and its denoised BEV representations are transferred to BEV perception
models through masked generative distillation~\citep{MGD}. At inference time,
the diffusion model and semantic branch are removed, introducing no additional
inference cost.


Our main contributions are as follows:
\begin{itemize}
    \item We construct a task-agnostic dense semantic BEV prior by projecting
    semantic masks extracted by a vision foundation model into the BEV space and
    introduce SBAM to inject this prior into diffusion features, allowing dense
    structures of both dynamic and static scene elements to guide denoising.
    
    \item We decouple diffusion training from downstream task supervision
    through generative BEV pre-training using only a denoising objective. This
    enables the BEV diffusion model to learn structured representations without
    relying on task-specific supervision.
    
    \item  We use masked generative distillation to transfer the pre-trained
    semantic-guided diffusion representations to BEV perception models,
    enabling competitive performance against task-specific diffusion approaches
    while preserving the inference cost of the underlying BEV model.
\end{itemize}


\begin{figure}[t]
  \centering
  \includegraphics[width=\linewidth]{figures/fig6.pdf}
  \caption{Overview of the proposed diffusion-based BEV representation learning
  framework. A frozen semantic segmentation model extracts dense semantic masks
  from multi-view images, which are projected into a semantic BEV prior. The
  prior guides diffusion-based denoising of the BEV features produced by the
  image backbone and BEV encoder.}
  \label{fig:overview}
\end{figure}

\section{Related Work}
\textbf{Diffusion Models for BEV Perception}
Diffusion models have demonstrated strong performance in generative modeling and
have increasingly been applied to discriminative tasks, including object
detection and semantic segmentation~\citep{DDPM,LDM,DDet,Segdiff}. Leveraging
their ability to recover structured representations through iterative
denoising, recent studies have extended diffusion models to BEV
perception~\citep{DiffBEV,DifFUSER,DiffVecMap,BEVDiff,LaneDiff}. DiffBEV and
DifFUSER use learned intermediate features as denoising conditions, which may
inherit errors from the upstream perception pipeline~\citep{DiffBEV,DifFUSER}.
To provide more reliable guidance, BEVDiffuser uses a ground-truth object layout
constructed from object categories and 3D bounding boxes~\citep{BEVDiff},
whereas LaneDiffusion injects lane centerline priors into BEV
features~\citep{LaneDiff}. Although effective, these priors depend on
task-specific annotations or representations and dedicated injection mechanisms,
limiting their applicability across tasks. Existing approaches also optimize
diffusion models jointly with task-specific objectives, coupling diffusion
training to particular downstream
tasks~\citep{DiffBEV,DifFUSER,DiffVecMap,BEVDiff,LaneDiff}. Extending these
methods to a new task therefore requires adapting the downstream task head and
training objectives. In contrast, our method decouples diffusion pre-training
from task-specific supervision and relies solely on a denoising objective.

\textbf{Semantic Conditioning in Generative Models}
Semantic layouts, such as segmentation maps, provide dense, spatially aligned
guidance for controllable image generation. SPADE modulates normalized
activations using spatially-adaptive affine parameters predicted from semantic
masks~\citep{SPADE}. Extending this principle to diffusion models, SDM
separately processes noisy images and semantic layouts and injects the layout
information into multiple decoder layers through spatially adaptive
normalization~\citep{SDM}. Other approaches condition pre-trained text-to-image
diffusion models through additional trainable
modules~\citep{Controlnet,GLIGEN,T2IAdapter}. ControlNet and T2I-Adapter inject
external conditions through trainable branches and lightweight adapters,
respectively~\citep{Controlnet,T2IAdapter}. Uni-ControlNet unifies multiple
local conditions, including segmentation maps, within a single model and injects
their multi-scale features through Feature Denormalization (FDN), which
spatially modulates normalized diffusion features~\citep{Uni-C}. Inspired by
this line of work, we develop Semantic BEV Adaptive Modulation (SBAM) to inject 
a dense semantic BEV prior into BEV diffusion features. Unlike image synthesis methods, 
our semantic condition is constructed by projecting multi-view semantic masks into 
the BEV space and is used to learn task-agnostic scene representations rather than 
synthesize images.



\section{Method}
\subsection{Preliminaries: Diffusion Models}
Diffusion models define a forward corruption process that gradually perturbs a
clean sample with Gaussian noise and learn a reverse process that removes the
noise~\citep{DDPM,DDIM}. Let $\mathbf{x}_0$ denote a clean BEV feature and
$t \in \{1,\dots,T\}$ denote a diffusion timestep. The forward process samples a
noisy feature $\mathbf{x}_t$ as
\begin{equation}
\label{eq:forward-diffusion}
  \mathbf{x}_t =
  \sqrt{\bar{\alpha}_t}\mathbf{x}_0 +
  \sqrt{1-\bar{\alpha}_t}\boldsymbol{\epsilon},
  \qquad
  \boldsymbol{\epsilon} \sim \mathcal{N}(\mathbf{0}, \mathbf{I}),
\end{equation}
where $\bar{\alpha}_t$ is the cumulative product of the noise schedule. A
denoising network is then trained to recover information about
$\mathbf{x}_0$ from $\mathbf{x}_t$ and $t$.

Different parameterizations can be used for the denoising target, such as
predicting the added noise, the clean sample, or the velocity. In this work, we
use the sample-prediction parameterization and train the denoising network
$f_{\theta}$ to predict the clean BEV feature. Given a condition $\mathbf{c}$,
the generic diffusion objective is
\begin{equation}
\label{eq:diffusion-objective}
  \mathcal{L}_{\mathrm{diff}} =
  \mathbb{E}_{\mathbf{x}_0,t,\boldsymbol{\epsilon}}
  \left[
  \left\|f_{\theta}(\mathbf{x}_t,t,\mathbf{c}) - \mathbf{x}_0\right\|_2^2
  \right].
\end{equation}
In the following sections, the condition $\mathbf{c}$ is instantiated as a
dense semantic BEV prior.



\subsection{Semantic-Guided BEV Diffusion Pre-training}

We propose a semantic-guided BEV diffusion framework that uses dense semantic
guidance in place of task-specific conditioning. As illustrated in
Figure~\ref{fig:overview}, multi-view semantic masks are projected into the BEV
space to form a dense semantic BEV prior for BEV feature denoising. Training
proceeds in two stages: generative diffusion pre-training without downstream
task supervision, followed by feature distillation from the frozen diffusion
model to downstream BEV perception models. At inference time, the diffusion
model and semantic branch are not used, preserving the inference cost of the
underlying BEV model.








\section{Experiments}
% Describe your experiments and results.

\section{Conclusion}
% Summarize your findings and future work.


% \begin{table}[t]
% \caption{Sample table title}
% \label{sample-table}
% \begin{center}
% \begin{tabular}{ll}
% \multicolumn{1}{c}{\bf PART}  &\multicolumn{1}{c}{\bf DESCRIPTION}
% \\ \hline \\
% Dendrite         &Input terminal \\
% Axon             &Output terminal \\
% Soma             &Cell body (contains cell nucleus) \\
% \end{tabular}
% \end{center}
% \end{table}

\section{Default Notation}

In an attempt to encourage standardized notation, we have included the
notation file from the textbook, \textit{Deep Learning}
\cite{goodfellow2016deep} available at
\url{https://github.com/goodfeli/dlbook_notation/}.  Use of this style
is not required and can be disabled by commenting out
\texttt{math\_commands.tex}.


\centerline{\bf Numbers and Arrays}
\bgroup
\def\arraystretch{1.5}
\begin{tabular}{p{1in}p{3.25in}}
$\displaystyle a$ & A scalar (integer or real)\\
$\displaystyle \va$ & A vector\\
$\displaystyle \mA$ & A matrix\\
$\displaystyle \tA$ & A tensor\\
$\displaystyle \mI_n$ & Identity matrix with $n$ rows and $n$ columns\\
$\displaystyle \mI$ & Identity matrix with dimensionality implied by context\\
$\displaystyle \ve^{(i)}$ & Standard basis vector $[0,\dots,0,1,0,\dots,0]$ with a 1 at position $i$\\
$\displaystyle \text{diag}(\va)$ & A square, diagonal matrix with diagonal entries given by $\va$\\
$\displaystyle \ra$ & A scalar random variable\\
$\displaystyle \rva$ & A vector-valued random variable\\
$\displaystyle \rmA$ & A matrix-valued random variable\\
\end{tabular}
\egroup
\vspace{0.25cm}

\centerline{\bf Sets and Graphs}
\bgroup
\def\arraystretch{1.5}

\begin{tabular}{p{1.25in}p{3.25in}}
$\displaystyle \sA$ & A set\\
$\displaystyle \R$ & The set of real numbers \\
$\displaystyle \{0, 1\}$ & The set containing 0 and 1 \\
$\displaystyle \{0, 1, \dots, n \}$ & The set of all integers between $0$ and $n$\\
$\displaystyle [a, b]$ & The real interval including $a$ and $b$\\
$\displaystyle (a, b]$ & The real interval excluding $a$ but including $b$\\
$\displaystyle \sA \backslash \sB$ & Set subtraction, i.e., the set
containing the elements of $\sA$ that are not in $\sB$\\
$\displaystyle \gG$ & A graph\\
$\displaystyle \parents_\gG(\ervx_i)$ & The parents of $\ervx_i$ in $\gG$
\end{tabular}
\vspace{0.25cm}


\centerline{\bf Indexing}
\bgroup
\def\arraystretch{1.5}

\begin{tabular}{p{1.25in}p{3.25in}}
$\displaystyle \eva_i$ & Element $i$ of vector $\va$, with indexing starting at 1 \\
$\displaystyle \eva_{-i}$ & All elements of vector $\va$ except for element $i$ \\
$\displaystyle \emA_{i,j}$ & Element $i, j$ of matrix $\mA$ \\
$\displaystyle \mA_{i, :}$ & Row $i$ of matrix $\mA$ \\
$\displaystyle \mA_{:, i}$ & Column $i$ of matrix $\mA$ \\
$\displaystyle \etA_{i, j, k}$ & Element $(i, j, k)$ of a 3-D tensor $\tA$\\
$\displaystyle \tA_{:, :, i}$ & 2-D slice of a 3-D tensor\\
$\displaystyle \erva_i$ & Element $i$ of the random vector $\rva$ \\
\end{tabular}
\egroup
\vspace{0.25cm}


\centerline{\bf Calculus}
\bgroup
\def\arraystretch{1.5}
\begin{tabular}{p{1.25in}p{3.25in}}
% NOTE: the [2ex] on the next line adds extra height to that row of the table.
% Without that command, the fraction on the first line is too tall and collides
% with the fraction on the second line.
$\displaystyle\frac{d y} {d x}$ & Derivative of $y$ with respect to $x$\\ [2ex]
$\displaystyle \frac{\partial y} {\partial x} $ & Partial derivative of $y$ with respect to $x$ \\
$\displaystyle \nabla_\vx y $ & Gradient of $y$ with respect to $\vx$ \\
$\displaystyle \nabla_\mX y $ & Matrix derivatives of $y$ with respect to $\mX$ \\
$\displaystyle \nabla_\tX y $ & Tensor containing derivatives of $y$ with respect to $\tX$ \\
$\displaystyle \frac{\partial f}{\partial \vx} $ & Jacobian matrix
$\mJ \in \R^{m\times n}$ of $f: \R^n \rightarrow \R^m$\\
$\displaystyle \nabla_\vx^2 f(\vx)\text{ or }\mH( f)(\vx)$ & The Hessian matrix of $f$ at input point $\vx$\\
$\displaystyle \int f(\vx) d\vx $ & Definite integral over the entire domain of $\vx$ \\
$\displaystyle \int_\sS f(\vx) d\vx$ & Definite integral with respect to $\vx$ over the set $\sS$ \\
\end{tabular}
\egroup
\vspace{0.25cm}

\centerline{\bf Probability and Information Theory}
\bgroup
\def\arraystretch{1.5}
\begin{tabular}{p{1.25in}p{3.25in}}
$\displaystyle P(\ra)$ & A probability distribution over a discrete variable\\
$\displaystyle p(\ra)$ & A probability distribution over a continuous variable, or over
a variable whose type has not been specified\\
$\displaystyle \ra \sim P$ & Random variable $\ra$ has distribution $P$\\
% so thing on left of \sim should always be a random variable, with name beginning with \r
$\displaystyle  \E_{\rx\sim P} [ f(x) ]\text{ or } \E f(x)$ & Expectation of $f(x)$ with respect to $P(\rx)$ \\
$\displaystyle \Var(f(x)) $ &  Variance of $f(x)$ under $P(\rx)$ \\
$\displaystyle \Cov(f(x),g(x)) $ & Covariance of $f(x)$ and $g(x)$ under $P(\rx)$\\
$\displaystyle H(\rx) $ & Shannon entropy of the random variable $\rx$\\
$\displaystyle \KL ( P \Vert Q ) $ & Kullback-Leibler divergence of P and Q \\
$\displaystyle \mathcal{N} ( \vx ; \vmu , \mSigma)$ & Gaussian distribution %
over $\vx$ with mean $\vmu$ and covariance $\mSigma$ \\
\end{tabular}
\egroup
\vspace{0.25cm}

\centerline{\bf Functions}
\bgroup
\def\arraystretch{1.5}
\begin{tabular}{p{1.25in}p{3.25in}}
$\displaystyle f: \sA \rightarrow \sB$ & The function $f$ with domain $\sA$ and range $\sB$\\
$\displaystyle f \circ g $ & Composition of the functions $f$ and $g$ \\
  $\displaystyle f(\vx ; \vtheta) $ & A function of $\vx$ parametrized by $\vtheta$.
  (Sometimes we write $f(\vx)$ and omit the argument $\vtheta$ to lighten notation) \\
$\displaystyle \log x$ & Natural logarithm of $x$ \\
$\displaystyle \sigma(x)$ & Logistic sigmoid, $\displaystyle \frac{1} {1 + \exp(-x)}$ \\
$\displaystyle \zeta(x)$ & Softplus, $\log(1 + \exp(x))$ \\
$\displaystyle || \vx ||_p $ & $\normlp$ norm of $\vx$ \\
$\displaystyle || \vx || $ & $\normltwo$ norm of $\vx$ \\
$\displaystyle x^+$ & Positive part of $x$, i.e., $\max(0,x)$\\
$\displaystyle \1_\mathrm{condition}$ & is 1 if the condition is true, 0 otherwise\\
\end{tabular}
\egroup
\vspace{0.25cm}



\section{Final instructions}
Do not change any aspects of the formatting parameters in the style files.
In particular, do not modify the width or length of the rectangle the text
should fit into, and do not change font sizes (except perhaps in the
\textsc{References} section; see below). Please note that pages should be
numbered.

\section{Preparing PostScript or PDF files}

Please prepare PostScript or PDF files with paper size ``US Letter'', and
not, for example, ``A4''. The -t
letter option on dvips will produce US Letter files.

Consider directly generating PDF files using \verb+pdflatex+
(especially if you are a MiKTeX user).
PDF figures must be substituted for EPS figures, however.

Otherwise, please generate your PostScript and PDF files with the following commands:
\begin{verbatim}
dvips mypaper.dvi -t letter -Ppdf -G0 -o mypaper.ps
ps2pdf mypaper.ps mypaper.pdf
\end{verbatim}

\subsection{Margins in LaTeX}

Most of the margin problems come from figures positioned by hand using
\verb+\special+ or other commands. We suggest using the command
\verb+\includegraphics+
from the graphicx package. Always specify the figure width as a multiple of
the line width as in the example below using .eps graphics
\begin{verbatim}
   \usepackage[dvips]{graphicx} ...
   \includegraphics[width=0.8\linewidth]{myfile.eps}
\end{verbatim}
or % Apr 2009 addition
\begin{verbatim}
   \usepackage[pdftex]{graphicx} ...
   \includegraphics[width=0.8\linewidth]{myfile.pdf}
\end{verbatim}
for .pdf graphics.
See section~4.4 in the graphics bundle documentation
(\url{http://www.ctan.org/tex-archive/macros/latex/required/graphics/grfguide.ps})


\subsubsection*{Author Contributions}
If you'd like to, you may include  a section for author contributions as is done
in many journals. This is optional and at the discretion of the authors.

\subsubsection*{Acknowledgments}
Use unnumbered third level headings for the acknowledgments. All
acknowledgments, including those to funding agencies, go at the end of the paper.


\bibliography{iclr2027_conference}
\bibliographystyle{iclr2027_conference}

\appendix
\section{Appendix}
You may include other additional sections here.


\end{document}
